\documentclass{article}
% \documentclass{amsart}

\usepackage{amsfonts}
\usepackage{amsmath}
\usepackage{biblatex}
\usepackage{blindtext}
\usepackage{geometry}
\usepackage{mathtools}
\usepackage{multicol}
\usepackage{xcolor}

\addbibresource{bibliography.bib}
\emergencystretch=1em %For the bibliography warning

% TODO:Change size appropriately
\geometry{a4paper, total={7in, 9.1in}} 

\begin{document}

\newcommand{\R}{\mathbb{R}}
\newcommand\tbf[1]{\textbf{#1}}
\newcommand\myworries[1]{\textcolor{red}{\tbf{#1}}}
% \renewcommand\myworries[1]{}

\title{Part B title}
\author{Fernando Herrera}
\date{\today}
\maketitle

\begin{multicols}{2}
\section{Introduction}
TODO \cite{thompson12}
\section{Core idea}
The goal of this section is to exemplify the technique used to obtain the electromagnetic cloak in a simpler setting.
To this end we will use the Helmholtz equation:
\begin{equation}
	\label{helmholtz}
	(\Delta + k^2)f=(\partial_x^2+\partial_y^2 + k^2)f=0
\end{equation}

And the change of coordinates
\begin{equation}
\begin{split}
	u=3x\\
	v=3y
\end{split}
\end{equation}
So we can think of this as a stretching of the space by a factor of 3.

How does equation \eqref{helmholtz} look like on the \(u,v\) coordinates?

Using the chain rule we have
\begin{equation}
\begin{split}
	\partial_x=\frac{du}{dx}\partial_u=3\partial_u\\
	\partial_y=\frac{du}{dy}\partial_v=3\partial_v
\end{split}
\end{equation}
so
\begin{equation}
\begin{split}
	\partial_x^2=\frac{du}{dx}\partial_u=9\partial_u^2\\
	\partial_y^2=\frac{du}{dy}\partial_v=9\partial_v^2
\end{split}
\end{equation}
Replacing this in \eqref{helmholtz} we get:
\begin{equation}
(9\partial_u^2+9\partial_v^2 +k^2)f=0\\
\end{equation}
and dividing by 9 the equation becomes:
\begin{equation}
(\partial_u^2+\partial_v^2 +(k/3)^2)f=0,
\end{equation}
which is exactly \eqref{helmholtz} with \(k\) replaced by \(k/3\).
This means that solving \eqref{helmholtz} with \(k/3\) is equivalent to
stretching the space by a factor of 3.

\end{multicols}
\printbibliography

\end{document}
